\section{Deep Q-Network (DQN)}
\label{sec:dqn}

\subsection{Algorithm Overview}

Deep Q-Networks (DQN)~\cite{mnih2015humanlevel} extend Q-learning by approximating the action-value function \(Q(s,a)\) with a neural network. The objective is to learn a policy that maximizes the expected discounted return by minimizing the temporal-difference (TD) error derived from the Bellman optimality equation:
\begin{equation}
    Q^*(s,a) = \mathbb{E}\left[r + \gamma \max_{a'} Q^*(s', a')\right].
\end{equation}

DQN optimizes the squared TD loss:
\begin{equation}
    \mathcal{L}(\theta) = \mathbb{E}_{(s,a,r,s') \sim \mathcal{D}} \left[ {\left( Q_\theta(s,a) - y \right)}^2 \right],
\end{equation}
where the target is defined as:
\begin{equation}
    y = r + \gamma \max_{a'} Q_{\theta^-}(s', a').
\end{equation}
Here, \(\theta^-\) are the parameters of a target network, updated periodically to stabilize training.

To improve data efficiency and reduce temporal correlations, transitions are stored in a replay buffer and sampled uniformly during training. Exploration is handled via an \(\epsilon\)-greedy policy, where random actions are selected with probability \(\epsilon\), which is annealed over time.

\subsection{Handling Discretized Continuous Control}

Applying DQN to continuous control requires discretizing the action space. In standard DQN, this discretization implies a combinatorial explosion of $K^N$ possible actions for an environment with $N$ independent joints and $K$ bins per joint. 

To bypass this exponential growth, we employ a Branching Dueling Q-Network (BDQN)~\cite{tavakoli2019action}. BDQN uses a shared state representation that branches into $N$ independent action-value estimators, reducing network outputs from $K^N$ to $K \times N$. This implicitly decouples joint actuation and it drastically saves parameters. Empirical tests on the Hopper environment showed BDQN achieves similar results to standard combinatorial DQN. Given its superior parameter efficiency, we adopt BDQN as our default architecture to scale effectively to more complex environments.

Additionally, default discretization wrappers map continuous actions to bin centers. We observed that at low resolutions (e.g., $K=3$), this artificially restricts the agent's torque range, limiting Hopper actions to approximately $[-0.6, 0.6]$ rather than the environment's full $[-1, 1]$ bounds, which severely impedes learning. To solve this, we use a custom discretization that maps to interval extremes, ensuring absolute maximum and minimum forces remain accessible.

\subsection{Custom Reward Shaping}
\label{sec:reward_shaping}

Standard implementations of continuous control environments often rely on a linear combination of proportional forward velocity, a constant healthy bonus, and a control penalty based on the absolute magnitude of the current action. To promote smoother gaits and more stable learning from high-dimensional pixel inputs, we heavily modified the reward function. 

First, rather than incentivizing unbounded forward acceleration, we encourage the agent to track a specific target velocity. This is achieved using an exponential kernel that smoothly penalizes deviations from the target speed:
\begin{equation}
    r_{\text{vel}} = w_{\text{vel}} \exp\left( - (v_x - v_{\text{target}})^2 \right),
\end{equation}
where $v_x$ is the current forward velocity, $v_{\text{target}}$ is the desired speed, and $w_{\text{vel}}$ is a forward weight. Second, we replace the standard absolute-value control penalty with a temporal smoothing cost. To actively penalize erratic, jittery movements, the control cost is computed as the squared difference between consecutive actions:
\begin{equation}
    r_{\text{ctrl}} = -w_{\text{ctrl}} \sum_{i} (a_{t, i} - a_{t-1, i})^2,
\end{equation}
with a penalty weight $w_{\text{ctrl}}$. The healthy survival reward is retained with a constant value $w_{\text{healthy}}$.

% --- FLOATING TOP PLOT FOR HOPPER ---
\begin{figure*}[t]
    \centering
    \includegraphics[width=\textwidth]{hopper_rainbow_reward_speed_x.pdf}
    \caption{Learning dynamics of Rainbow on the Hopper environment. The agent rapidly converges to a stable, high-speed gait in approximately 10M steps.}
    \label{fig:hopper_rainbow_results}
\end{figure*}
% -----------------------------------

\subsection{Hopper: Learning from Pixels}
\label{sec:hopper_results}

Before scaling to more complex challenges, we validate our algorithm on a simpler continuous control task: the Hopper environment. The primary objective is to verify that the agent can learn stable locomotion relying entirely on visual inputs rather than standard low-dimensional state vectors. 

We found the following configurations to be crucial for successful training:

\textbf{Environment Configuration:} 
The agent operates directly on pixel data, observing $64 \times 64$ images with a frame stack of $4$ to capture temporal dynamics. We deploy $30$ parallel environments to increase data throughput and diversify the replay buffer. Notably, we widen the \texttt{healthy\_angle\_range} to $[-0.4, 0.4]$ (from the default $[-0.2, 0.2]$), allowing the agent to explore more diverse states during the early stages of training.

\textbf{Reward Configuration:}
The reward function is shaped to encourage a target velocity of $1.5$ units/s, with a velocity reward weight of $w_{\text{vel}} = 1.0$ and weight of $w_{\text{vel}} = 3.0$. The control penalty weight is set to $w_{\text{ctrl}} = 0.001$ to promote smooth actions, while the healthy survival bonus is maintained at $w_{\text{healthy}} = 1.0$.

\textbf{Algorithm Configuration:} 
We train the Branching DQN (BDQN) for $50,000,000$ environment frames. The action space is mapped to $3$ discrete bins representing the extremes of the joint torque intervals. To handle the instability and large Temporal-Difference (TD) errors common when training Q-networks from pixels, we utilize a Huber loss function rather than standard MSE. Exploration follows an $\epsilon$-greedy policy, annealed linearly over a prolonged period of $5,000,000$ frames. Optimization is performed using a learning rate of $10^{-4}$, with the target network updated every $10,000$ frames to provide stable bootstrapping targets.

The learning dynamics of this configuration are illustrated in Figure~\ref{fig:hopper_results}. As training progresses, the agent's step-wise speed smoothly converges toward the target velocity of $1.5$. Note that it converges to $1.25$, which indicates that for the strategy it learned that is the most stable velocity. A higher target may have learned a faster strategy or just collapsed. This convergence is accompanied by a steady increase in both the episodic reward and the final $X$ position, indicating that the combination of our custom reward shaping and the BDQN architecture successfully enables the agent tomaintain balance and move forward relying strictly on high-dimensional pixel observations.

% --- SINGLE COLUMN PLOT PLACEHOLDER ---
\begin{figure}[htbp]
    \centering
    \includegraphics[width=\linewidth]{dqn_rewards.pdf}
    \caption{Training reward across multiple hyperparameter sweeps on the Humanoid environment using DQN.}
    \label{fig:humanoid_failures}
\end{figure}
% ------------------------------------

\subsection{Humanoid: The Limits of Discretized DQN}
\label{sec:humanoid_results}

Following the successful application of Branching DQN on the Hopper environment, we attempted to scale the algorithm to the highly complex Humanoid environment. This presents a massive leap in difficulty due to its inherent bipedal instability and a vast continuous action space comprising 17 independent joints. 

\textbf{Environment Tweaks:}
To mitigate this complexity, we artificially constrained the morphology by freezing several non-essential actuators. Specifically, we locked the Z-axis rotations of the abdomen and hips, and completely disabled all six arm joints (rendering the arms physically limp). By freezing these 9 specific joints ($[1, 4, 8, 11, 12, 13, 14, 15, 16]$), we force the agent to focus purely on the core mechanics of leg actuation and balance. Additionally, to capture the finer nuances of this complex body in the visual state, we slightly increased the observation resolution from $64 \times 64$ to $96 \times 96$ pixels.

\textbf{Hyperparameter Sweeps:}
To find a viable learning configuration, we conducted an exhaustive hyperparameter search. We massively scaled the training horizon (up to 100M frames) and prolonged the $\epsilon$-greedy exploration phase (annealing over up to 50M frames). We tuned replay buffer capacities (up to 500k), learning rates ($10^{-5}$ to $10^{-4}$), batch sizes, and target network update frequencies. Furthermore, we experimented heavily with the reward shaping parameters, tweaking target velocities, contact costs, and survival bonuses to guide the agent.

Despite these extensive morphological simplifications and exhaustive hyperparameter sweeps, none of the trials successfully learned to stably walk (see Figure~\ref{fig:humanoid_failures}). These negative results strongly indicate the fundamental limitations of Q-learning and action discretization—even when mitigated by branching architectures—for high-dimensional continuous control tasks where joint actuations are deeply correlated and highly sensitive to approximation errors.